\section{Proof of the main theorem}\label{sec:completion}
We now combine the construction with the rounding estimate. The order of
the choices is important: the accuracy parameter fixes all auxiliary data,
and only then does the length tend to infinity.

\begin{proof}[Proof of Theorem~\ref{thm:main}]
Let $b=1/8$ and let $D>0$ be the absolute constant in
Proposition~\ref{prop:correction}.
For a fixed sufficiently small $\delta>0$, take the real coefficients
$Y_{N,k}$ and the functions $B_N$ from Proposition~\ref{prop:starting},
and the correction $R_N$ from Proposition~\ref{prop:correction}.
All estimates below hold for every sufficiently large integer $N$.
Put
\[
 s_\delta=D\sqrt\delta,\qquad S_\delta=1+s_\delta,
 \qquad
 X_{N,k}=\frac{Y_{N,k}+\sqrt N\,\widehat R_N(k)}{S_\delta}
 \quad(0\le k<N).
\]
Conjugate symmetry of $R_N$ makes $\widehat R_N(k)$ real. Since
$|Y_{N,k}|\le1$ and $\sqrt N|\widehat R_N(k)|\le s_\delta$,
the vector $X_N$ belongs to $[-1,1]^N$.
The uniform approximation in Proposition~\ref{prop:starting} and the
exterior Fourier-tail estimate in Proposition~\ref{prop:correction} give
\begin{equation}\label{eq:completion-profile}
 U_{X_N}(t)=\frac{B_N(t)+R_N(t)}{S_\delta}+o(1)
 \qquad\text{uniformly for }t\in\T.
\end{equation}
Consequently,
\begin{equation}\label{eq:completion-envelope}
 \min_{t\in\T}|U_{X_N}(t)|\ge\frac b{S_\delta}-o(1),
 \qquad
 \norm{U_{X_N}}_\infty\le K_\delta+o(1).
\end{equation}
In the second inequality we used $S_\delta\ge1$.

The same coefficient bound controls the cost of rounding. Directly,
\[
 |X_{N,k}-Y_{N,k}|
 =\frac{|\sqrt N\,\widehat R_N(k)-s_\delta Y_{N,k}|}{S_\delta}
 \le2s_\delta.
\]
Both coefficients lie in $[-1,1]$, so
$|X_{N,k}^2-Y_{N,k}^2|\le4s_\delta$.
Using the mean-square estimate from Proposition~\ref{prop:starting} and
$1-|x|\le1-x^2$ for $|x|\le1$, we obtain
\begin{equation}\label{eq:completion-defect}
 \frac{\mu(X_N)}N
 \le\frac1{2N}\sum_{k=0}^{N-1}(1-X_{N,k}^2)
 \le4\delta+2D\sqrt\delta=:q_\delta.
\end{equation}
Choose $\delta$ small enough that $q_\delta<1/2$.
The function $q\mapsto q\log(80/q)$ is increasing on $(0,1/2]$.
Lemma~\ref{lem:rounding} therefore supplies signs
$\varepsilon_0,\ldots,\varepsilon_{N-1}$ with
\begin{equation}\label{eq:completion-rounding}
 \norm{U_\varepsilon-U_{X_N}}_\infty
 \le \frac C{\sqrt N}
       +C\sqrt{q_\delta\log(80/q_\delta)}.
\end{equation}
If $\mu(X_N)=0$, we instead take $\varepsilon_k=X_{N,k}$, so this bound
also covers that case.

Fix $c=1/16$ once and for all. As $\delta\downarrow0$,
\[
 \frac b{S_\delta}
       -C\sqrt{q_\delta\log(80/q_\delta)}\longrightarrow\frac18,
 \qquad
 K_\delta+C\sqrt{q_\delta\log(80/q_\delta)}\longrightarrow1.
\]
For any prescribed $\eta>0$, first choose $\delta$ so small that the
first expression is strictly larger than $c$ and the second is strictly
smaller than $1+\eta$. Fix every construction datum associated with this
$\delta$. We may then choose $N_0$ so that for every $N\ge N_0$ all preceding
estimates hold and the errors in \eqref{eq:completion-envelope} and
\eqref{eq:completion-rounding} are smaller than these strict margins.
It follows that
\[
 c\le|U_\varepsilon(t)|\le1+\eta\qquad(t\in\T).
\]
Multiplication by $\sqrt N$ gives the asserted polynomial. Rounding assigns
a real sign to each of the consecutive indices $0,\ldots,N-1$; all
estimates include $t=0$ and $t=1/2$ and hold through every sufficiently
large integer length.
\end{proof}
